\section{Evaluación y resultados}
A continuación se resume cómo evaluamos UrbanTracker y qué observamos en las pruebas. Nos centramos en tres cosas: que las funciones clave se comporten bien (funcional), que el tiempo real sea usable y que el sistema se recupere si hay problemas de red (resiliencia). Las decisiones técnicas (API REST documentada, canal ligero para eventos y separación por capas) van en la misma línea de lo que reporta la literatura en transporte y sistemas híbridos \cite{cityruta2018new}; \cite{stopbus2016new}; \cite{sismo2016}; \cite{easyrest2022}; \cite{ingenieria2022}.

\subsection{Diseño de la evaluación}
Funcional: login de conductor/administrador, elección de ruta, inicio/fin de servicio, CRUD en el panel y visualización en el mapa.

Tiempo real: frecuencia de envío desde la app del conductor y latencia extremo a extremo (del teléfono a la pantalla web).

Resiliencia: qué pasa si se pierde la señal, si se revocan permisos de ubicación o si hay reconexiones.
La separación por capas y tener la API documentada ayudó a reducir fricción al integrar web, móvil y backend \cite{easyrest2022}, y la organización del código siguiendo principios de POO favoreció cambios sin romper lo anterior \cite{poo2023}.

\subsection{Métricas observadas}
Frecuencia de actualización: del orden de ~5 s entre envíos de posición desde la app del conductor.

Latencia extremo a extremo: < 2 s desde que el móvil envía hasta que el usuario lo ve en la web (en 4G estable).

Concurrencia inicial: 50+ clientes suscritos sin degradación visible.

Estos resultados son coherentes con soluciones cercanas en transporte/logística que manejan grandes volúmenes de eventos en vivo \cite{torres2020}; \cite{odesanya2025}; \cite{anyango2025} y con el uso de MQTT como mensajería ligera en contextos vehiculares \cite{sismo2016}.

\subsection{Evolución y resiliencia}
En sesiones repetidas el mapa se mantuvo estable y, cuando hubo cortes de red, la app del conductor reintentó y siguió enviando posiciones al recuperar conectividad. Mantener el contrato de API claro y separar responsabilidades facilitó depurar problemas y repetir pruebas sin sorpresas \cite{easyrest2022}; \cite{poo2023}. De cara al crecimiento, las guías de sistemas híbridos sirven para decidir si agregar almacenamiento específico cuando se quiera guardar trayectorias históricas \cite{ingenieria2022}.

\subsection{Interpretación de resultados}
Los valores observados (frecuencia ~5 s, latencia < 2 s y >50 suscriptores) son suficientes para la experiencia que queremos: el pasajero confía en lo que ve y el administrador puede tomar decisiones en el momento. Igual, el desempeño depende del contexto de red y de dónde se despliega el backend/broker \cite{cityruta2018new}; \cite{stopbus2016new}; \cite{ingenieria2022}. Como siguiente paso, tiene sentido instrumentar métricas continuas (series de tiempo) y, si se requiere análisis profundo, evaluar un componente adicional para histórico de eventos.